\section{Background and Related Work}
\label{sec:rl}

\noindent
\textbf{Continual Learning}: 
Continual learning approaches can be organized into a few broad categories which are all useful depending on the problem setting and constraints. One group of approaches expands a model's architecture as new tasks are encountered; these are highly effective for applications where a model growing with tasks is practical~\cite{ebrahimi2020adversarial,lee2020neural,lomonaco2017core50,maltoni2019continuous,Rusu:2016}. 
Another approach is to regularize the model with respect to past task knowledge while training the new task. This can either be done by regularizing the model in the  weight space (i.e., penalize changes to model parameters)~\cite{aljundi2017memory,ebrahimi2019uncertainty,kirkpatrick2017overcoming,titsias2019functional,zenke2017continual} or the prediction space (i.e., penalize changes to model predictions)~\cite{Ahn2021ssil,castro2018end,hou2018lifelong,lee2019overcoming,li2016learning}. Regularizing knowledge in the prediction space is done using \textit{knowledge distillation}~\cite{hinton2015distilling} and it has been found to perform better than model regularization-based methods for continual learning when task labels are not given~\cite{lesort2019generative,van2018generative}.

Rehearsal with stored data~\cite{aljundi2019online, aljundi2019gradient,bang2021rainbow,buzzega2020dark,chaudhry2018efficient,chaudhry2019episodic,Gepperth:2017,hayes2018memory,hou2019learning,Kemker:2017,Lopez-Paz:2017,Rebuffi:2016,robins1995catastrophic, rolnick2019experience,von2019continual,wu2019large} or samples from a generative model~\cite{kamra2017deep,kemker2018fear,ostapenko2019learning,Shin:2017,van2020brain} is highly effective when storing training data or training/saving a generative model is possible. Unfortunately, for many machine learning applications long-term storage of training data will violate data privacy, as well as incur a large memory cost. With respect to the generative model, this training process is much more computationally and memory intensive compared to a classification model and additionally may violate data legality concerns because using a generative model increases the chance of memorizing potentially sensitive data~\cite{nagarajan2018theoretical}. This motivates us to work on the important setting of \emph{rehearsal-free} approaches to mitigate catastrophic forgetting.

\noindent
\textbf{Rehearsal-Free Continual Learning}: 
Recent works propose producing images for rehearsal using deep-model inversion~\cite{choi2021dual,gao2022r,smith2021abd,yin2020dreaming}. While these methods perform well compared to generative modeling approaches and simply rehearse from a small number of stored images, model-inversion is a slow process associated with high computational costs in the continual learning setting, and furthermore these methods under-perform rehearsal-based approaches by significant margins~\cite{smith2021abd}. Other works have looked at rehearsal-free continual learning from an online ``streaming" learning perspective using a frozen, pre-trained model~\cite{hayes:2019,lomonaco2020rehearsal}. Because we allow our models to train to convergence on task data (as is common for continual learning~\cite{wu2019large}), our setting is very different. %

\noindent
\textbf{Continual Learning in Vision Transformers}:
Recent work~\cite{brown2020language,li2022technical,yin2019benchmarking} has proven transformers to generalize well to unseen domains. 
For example, one study varies the number of attention heads in Vision transformers and conclude that ViTs offer more robustness to forgetting with respect to CNN-based equivalents~\cite{mirzadeh2022architecture}. Another~\cite{yu2021improving} shows that the vanilla counterparts of ViTs are more prone to forgetting when trained from scratch. Finally, DyTox~\cite{douillard2022dytox} learns a unified model with a parameter-isolation approach and dynamically expands the tokens processed by the last layer to mitigate forgetting. For each task, they learn a new task-specific token per head using task-attention based decoder blocks. However, the above works either rely on exemplars to defy forgetting or they need to train a new transformer model from scratch, a costly operation, hence differing from our objective of exemplar-free CL using pre-trained ViTs.

\noindent
\textbf{Prompting for Continual Learning}: 
Prompt-based approaches for continual learning boast a strong protection against catastrophic forgetting by learning a small number of insert-able model instructions (prompts) rather than modifying encoder parameters directly. The current state-of-the-art approaches
for our setting, DualPrompt~\cite{wang2022dualprompt} and L2P~\cite{wang2022learning}, create a pool of prompts to be selected for insertion into the model, matching input data to prompts \emph{without task id} with a local \emph{clustering-like} optimization. We build upon the foundation of these methods, as discussed later in this paper. We note that the recent S-Prompts~\cite{wang2022sprompt} method is similar to these approaches but designed for \emph{domain-incremental learning} (i.e., learning the same set of classes under covariate distribution task shifts), which is different than the class-incremental continual learning setting of our paper (i.e., learn emerging object classes in new tasks).

